% =============================================================================
% test-abnt-projeto-pesquisa.tex
% Documento de teste para abnt-projeto-pesquisa.cls
%
% Demonstra um projeto de pesquisa conforme NBR 15287:2025, incluindo
% capa, folha de rosto, sumário, numeração progressiva, seções,
% alíneas e subalíneas.
%
% Compilação:
%   cd tests && TEXINPUTS=../src: latexmk -pdf -shell-escape -interaction=nonstopmode test-abnt-projeto-pesquisa.tex
%
% Checklist:
%   1  — Capa herdada de abnt.cls
%   2  — Folha de rosto herdada de abnt.cls
%   3  — Sumário carregado automaticamente (abnt-sumario)
%   4  — Numeração progressiva carregada automaticamente (abnt-numeracao)
%   5  — Seções primárias em maiúsculas e negrito
%   6  — Seções secundárias em maiúsculas sem negrito
%   7  — Seções terciárias em negrito
%   8  — Alíneas com letras minúsculas e parêntese
%   9  — Subalíneas com travessão
%   10 — Classe NÃO define \postextual, \imprimirdedicatoria, etc.
%   11 — Paginação correta: pré-textuais ocultos, textuais visíveis
% =============================================================================

\documentclass{abnt-projeto-pesquisa}

% Metadados
\titulo{Impacto da inteligência artificial na produtividade agrícola brasileira}
\autor{Lucas Gabriel Ferreira}
\orientador{Prof. Dr. Antônio Carlos Souza}
\instituicao{Universidade Federal de Goiás}
\local{Goiânia}
\ano{2025}
\natureza{Projeto de pesquisa apresentado ao Programa de
Pós-Graduação em Agronomia da Universidade Federal de Goiás.}
\programa{Programa de Pós-Graduação em Agronomia}
\area{Agricultura de Precisão}

\begin{document}

% =====================================================================
% PRÉ-TEXTUAIS (checklist 1, 2, 3, 11)
% =====================================================================
\pretextual

% Capa (checklist 1)
\imprimircapa

% Folha de rosto (checklist 2)
\imprimirfolhaderosto

% Sumário (checklist 3)
\imprimirsumario

% =====================================================================
% ELEMENTOS TEXTUAIS (checklist 4–9, 11)
% =====================================================================
\textual

\chapter{Introdução}

A agricultura brasileira enfrenta desafios crescentes relacionados à
produtividade e à sustentabilidade ambiental. A inteligência artificial
surge como ferramenta promissora para otimizar processos agrícolas.

\section{Contextualização}

O Brasil é um dos maiores produtores agrícolas do mundo. A adoção de
tecnologias digitais no campo tem se acelerado nos últimos anos.

\subsection{Agricultura de precisão}

A agricultura de precisão utiliza dados georreferenciados para a tomada
de decisão no manejo das culturas.

\chapter{Justificativa}

A pesquisa se justifica pela necessidade de avaliar o impacto real das
tecnologias de inteligência artificial na produtividade das lavouras
brasileiras.

\chapter{Objetivos}

\section{Objetivo geral}

Avaliar o impacto da inteligência artificial na produtividade agrícola
em propriedades rurais do cerrado brasileiro.

\section{Objetivos específicos}

Os objetivos específicos deste projeto são:

% Alíneas (checklist 8)
\begin{alineas}
  \item mapear as tecnologias de inteligência artificial adotadas
    nas propriedades selecionadas;
  \item comparar indicadores de produtividade antes e depois da
    adoção das tecnologias;
  \item identificar fatores que influenciam a adoção, incluindo:
    % Subalíneas (checklist 9)
    \begin{subalineas}
      \item custo de implementação;
      \item disponibilidade de conectividade;
      \item nível de capacitação técnica dos produtores.
    \end{subalineas}
\end{alineas}

\chapter{Metodologia}

A pesquisa adotará abordagem mista, combinando análise quantitativa de
dados de produtividade com entrevistas qualitativas junto aos produtores
rurais.

\section{Área de estudo}

O estudo será conduzido em propriedades rurais localizadas no estado de
Goiás, com foco na produção de soja e milho.

\section{Coleta de dados}

Os dados serão coletados por meio de sensores de campo, imagens de
satélite e questionários aplicados aos produtores.

\chapter{Cronograma}

O projeto será executado em vinte e quatro meses, divididos em quatro
etapas principais.

\chapter{Referências}

EMBRAPA. \textbf{Agricultura digital no Brasil}: tendências,
desafios e oportunidades. Brasília, 2023.

\end{document}
